% Generated draft: requires mathematical and source review.
\documentclass[11pt]{article}
\usepackage[margin=25mm]{geometry}
\usepackage{fontspec}
\setmainfont{Latin Modern Roman}
\usepackage{amsmath,amssymb,mathtools,unicode-math}
\setmathfont{Latin Modern Math}
\usepackage{enumitem,needspace,xurl,hyperref}
\setlength{\parindent}{0pt}
\setlength{\parskip}{5pt}
\setlength{\emergencystretch}{3em}
\newcommand{\N}{\mathbb{N}}
\newcommand{\Z}{\mathbb{Z}}
\newcommand{\Q}{\mathbb{Q}}
\newcommand{\R}{\mathbb{R}}
\newcommand{\C}{\mathbb{C}}
\newcommand{\F}{\mathbb{F}}
\newcommand{\A}{\mathbb{A}}
\newcommand{\PP}{\mathbb{P}}
\newcommand{\E}{\mathbb{E}}
\newcommand{\eps}{\varepsilon}
\newcommand{\dd}{\,\mathrm{d}}
\begin{document}
\section*{UnsolvedMath ID: 30005603}
\textbf{Source problem code:} catalogue entry
\subsection{Definitions and mathematical statement}
Let \(\mathbb H=\{x+iy:x,y\in\R,\ y>0\}\), and let \(\Gamma=\operatorname{PSL}_2(\Z)\) act by fractional linear transformations. A function on \(\Gamma\backslash\mathbb H\) is identified with its \(\Gamma\)-invariant lift to \(\mathbb H\). Such a lift is periodic in \(x\) with period \(1\). For a sufficiently regular lift \(f\), define its Fourier coefficients by
\[
 f_n(y)=\int_0^1 f(x+iy)e^{-2\pi i n x}\,\dd x,
 \qquad n\in\Z,\quad y>0.
\]
The corresponding Fourier expansion is \(f(x+iy)=\sum_{n\in\Z}f_n(y)e^{2\pi i n x}\), with its mode of convergence depending on the regularity assumptions.

The supplied record asks for all Fourier coefficients of solutions of
\[
 \bigl(\Delta-r(r+1)\bigr)f=E_aE_b
 \quad\text{on }\Gamma\backslash\mathbb H,
\]
and for evaluations or arithmetic information concerning the divisor-correlation sums occurring in those coefficients. Here \(\Delta\), \(E_a\), \(E_b\), and \(r\) retain the source's notation. The record does not specify the sign convention for the hyperbolic Laplacian, the definitions or normalization of the functions \(E_a,E_b\), or the allowed values of \(a,b,r\). Consequently, this equation records the requested problem but does not yet specify a complete analytic problem.

No regularity class, growth condition at the cusp, boundary condition, or normalization selecting \(f\) is supplied. In particular, whenever \(h\) is an admissible solution of \((\Delta-r(r+1))h=0\), the function \(f+h\) solves the same inhomogeneous equation. Thus the requested coefficients must either be described for all admissible solutions, including their homogeneous freedom, or be attached to a solution selected by additional conditions that are absent from the record.

For each positive integer \(n\), the original excerpt describes sums indexed by
\[
 I_n=\{(n_1,n_2)\in(\Z\setminus\{0\})^2:n_1+n_2=n\}.
\]
It visibly includes factors \(\sigma_0(n_1)\sigma_0(n_2)\) and a logarithmic expression. For a positive integer \(m\), the usual divisor-counting function is \(\sigma_0(m)=\sum_{d\mid m,\ d>0}1\); however, the source does not define its use at negative arguments. The extracted display is too corrupted to determine the full summand, including the logarithm's argument and the placement of powers and rational factors. The set \(I_n\) is infinite, and the record supplies no convergence or summation prescription. Therefore an exact divisor-correlation sum cannot be stated from this evidence. Recovering that display and its conventions is necessary before the arithmetic part has a precise formulation.
\subsection{Short English statement}
The supplied problem asks for the Fourier coefficients of solutions to an inhomogeneous automorphic Laplace equation on the modular surface, and for arithmetic evaluations of associated divisor-correlation sums. Parameter ranges, operator and Eisenstein-series conventions, and conditions selecting a solution are missing. The extracted divisor-sum formula is corrupted, so its exact evaluation target remains unresolved.
\subsection{Sources}
\begin{itemize}
\item Local registry: \texttt{data/unsolvedmath/problems.json}, permanent UnsolvedMath ID \(30005603\), source problem code \texttt{catalogue entry}, title ``Fourier Coefficients of Inhomogeneous Automorphic Laplace Equations.'' The supplied record gives proposed year \(2023\) and names no proposer.
\item Supplied primary-report citation: \emph{Geometric Spectral Theory (2024)}, Oberwolfach Reports, pp.~2017--2106, DOI \texttt{10.4171/owr/2023/36}. Supplied source URL and literature-source link, labeled ``Oberwolfach Report 36'': \url{https://doi.org/10.4171/owr/2023/36}. The background describes the collection as drawn from workshop reports and their open-problem discussions.
\item The supplied literature assessment, dated exactly \texttt{2026-08-22}, labels the problem open and states that the primary report leaves the Fourier-coefficient and divisor-sum questions open. This is supplied dated curation, not an independently verified literature assessment. Its accompanying qualification calls for independent verification of time-sensitive status and cited claims.
\item The registry labels the statement \texttt{corrected\_verified} and says it was reconstructed and checked against the official report. These are supplied provenance claims; this draft has not independently checked the report or verified the reconstruction.
\end{itemize}
\end{document}
